\Article{20}
Aus der Theorie der Elimination folgt leicht, dass die unvermeidlichen Beobachtungsfehler die Rechnung unsicherer machen. Je mehr Koeffizienten durch Elimination zu bestimmen sind, desto größer ist die Unsicherheit. Deshalb ist das in Artikel $18$, II. vorgeschriebene Verfahren besonders wertvoll, denn es unterdrückt die Koeffizienten der Potenzen $R^{-(n+2)},R^{-(n+4)}$. Bei vollkommener Symmetrie würden diese Terme zwar von selbst wegfallen, doch wäre es zu unsicher, sich darauf zu verlassen. Eine kleine Abweichung von der Symmetrie fiele bei der ersten Methode weit weniger ins Gewicht als bei der zweiten. Bei der ersten Methode achte man wenigstens auf den Punkt $h'$. Liegt dieser Punkt, von dem aus die Abstände gemessen werden, hinreichend genau im magnetischen Meridian durch $h$, so wird sich zwischen $u'-u''$ und $u'''-u''''$ kaum ein merklicher Unterschied zeigen. Bei der zweiten Methode verhält es sich anders, besonders wenn der Apparat eine exzentrische Aufhängung erfordert. Mit dieser Methode erzielt man stets eine viel geringere Genauigkeit, wenn der Raum keine Beobachtungen von beiden Seiten erlaubt. Die erste Methode ist noch aus einem weiteren Grund vorzuziehen. Da in der Natur $n=2$ ist, liefert sie einen doppelt so großen Wert von $L$ wie die zweite. Will man jedoch bei der zweiten Methode mit exzentrischer Aufhängung den von $R^{-(n+2)}$ abhängigen Term so weit wie möglich beseitigen, \SourceLeaf{113}so muss man den Punkt $h'$ so wählen, dass die Mitte der Nadel (für $u=u^0$) genau zwischen $h$ und $h'$ liegt. Die Rechnung, die dieses Ergebnis geliefert hat, muss hier jedoch der Kürze halber entfallen.

